% EXTRACTED FOR PAPER TWO (the hardware case study) -- PASS 48.1
% source     : paper/paper.tex
% blob       : 13590dbfe8119022ae95ff7a374910e30146ab8c
% lines      : 1094-1350 (1-indexed, inclusive)
% section    : Section 6, all six subsections
% prose words: 1691
%
% The hardware case study in full: platform and execution-path disclosure (6.1), the entangling-overhead measurement (6.2), the readout and SPAM characterization with the four-qubit table (6.3), cross-session variation with Figure 6 (6.4), the seven-mechanism elimination ledger as a longtable (6.5), and what remains (6.6). This is paper two's core.
%
% This is a VERBATIM COPY. The original is unmodified and still frozen.
% Regenerate with: experiments/pass48_extract_paper2.py

\hypertarget{hardware-a-pre-registered-test-on-rigetti-cepheus-1-108q}{%
\section{Hardware: a committed-in-advance test on Rigetti
Cepheus-1-108Q}\label{hardware-a-pre-registered-test-on-rigetti-cepheus-1-108q}}

\hypertarget{platform-and-a-disclosure-about-the-execution-path}{%
\subsection{Platform, and a disclosure about the execution
path}\label{platform-and-a-disclosure-about-the-execution-path}}

All measurements were submitted to the Rigetti Cepheus-1-108Q backend as
provisioned by the platform, a 108-qubit superconducting processor of
twelve interconnected nine-qubit chiplets~\cite{rigetti2026cepheus}, whose
published median fidelities at the time were 99.1\% for the two-qubit CZ
gate and 99.9\% for single-qubit gates. Readout error is absent from that
release, a fact that matters in Section 6.3. These are the specification
as published at that date; per-job calibration metadata was not available
through the execution path we used.

Access was through the Open Quantum platform operated by Quantum Rings
\cite{wold2025openquantum}, on the Public Tier, which does not route jobs
directly to the vendor: circuits go to distributed operators on the
Quantum Compute subnet (SN48) of the Bittensor network, with validators
spot-checking that they ran on the target QPU \cite{wold2025openquantum}.
We therefore cannot
independently verify the execution path of any individual job. This does
not affect the characterizations in Sections 6.2 and 6.3, which are
internally consistent, but it is a candidate contributor to the
session-to-session variability of Section 6.4, listed as such in Section
8. Jobs were submitted as OpenQASM 3 with physical-qubit addressing;
registers, protocol, license conditions, and the credit ledger are in the
supplementary material.

\hypertarget{the-entangling-overhead-is-nearly-free}{%
\subsection{The modeled entangling overhead is small for the tested circuits}\label{the-entangling-overhead-is-nearly-free}}

The destructive SWAP test at \(n = 2\) transpiles onto physical qubits
\(\{0, 1, 9, 10\}\) using four CZ gates and zero routing; the GHZ ladders
at \(n = 3, 4\) map with zero routing at \(3n - 2\) CZ gates and nest. CZ
error on these qubits was not measured directly: the platform compiler
resynthesized our characterization echo and optimized its gates away, and
the verbatim box that would have prevented this failed at execution, so
we assume the published median of 0.9\%, an average infidelity \(r\)
corresponding under \(\Lambda(\rho) = (1-p)\rho + p\,\mathbb{I}/d\) to
\(p = \tfrac{4}{3} r = 0.012\) for a two-qubit gate (\(d = 4\)).

On that assumption gate error accounts for roughly 0.042 of the purity
deficit at \(n = 2\) against a total deficit of 0.29, about a seventh
(roughly 14\%). At the lower median the vendor has since published for
this backend (\(98.71\%\), \(r = 0.0129\), \(p = 0.0172\)) it rises to
about 0.057, or close to a fifth of the measured deficit. This is a
device-wide median, not an edge-specific measurement, and a depolarizing
model captures neither coherent error nor leakage, so the figure is a
scale rather than a budget; under either median it is too small to account
for the deficit.

The resource-demand concern of \cite{qmladvantage2026} therefore does not
materialize for the destructive SWAP test on a square-lattice device, a
property of structured states on matched topology rather than of the test
in general, since generic states at \(n = 4\) require 46 CZ gates here, 26
intrinsic plus 20 routing-attributable, against the ladder's 10 with zero
routing, which is why we restrict the hardware series to GHZ ladders. This
is a \(k = 2\) statement; the ancilla-based Hadamard test for \(k \geq 3\)
(Section 2.3) carries a cost we have not measured.

\hypertarget{readout-is-the-wall-and-it-is-not-in-the-datasheet}{%
\subsection{Consistent with readout and SPAM error dominating, and not in the datasheet}\label{readout-is-the-wall-and-it-is-not-in-the-datasheet}}

Our locked prediction, from the published gate fidelities and an assumed
2\% readout error and committed to version control before submission, was
a Bell-state purity of 0.8932. The device returned

\[0.7184 \quad [0.6992,\, 0.7376] \;\; \text{(95\% bootstrap CI, 5000 shots)},\]

against a true purity of exactly 1, overshooting by more than seventeen
standard errors, conditional shot noise, hence a lower bound on the
total uncertainty (Appendix C). The estimand is the target the circuit was
written to prepare, so every deviation is charged to the
global-depolarizing channel of Section 4.1, which has no single-copy
reading to be relabelled into. The output structure is consistent with the
intended circuit: the four dominant outcomes are exactly the ideal support
of the Bell SWAP test, with noise leaking about 25\% of the weight into
the remaining twelve.

Inverting the bias law of Section 4.1 gives an effective global
depolarizing rate \(g = 0.375\), over two and a half times the
\(g = 0.142\) the locked prediction implied.

\begin{table}[htbp]
\centering
\setlength{\tabcolsep}{10pt}
\begin{tabular}{@{}lll@{}}
\toprule
qubit & \(P(1 \mid 0)\) & \(P(0 \mid 1)\) \\
\midrule
\(q_0\) & 9.3\% & 8.9\% \\
\(q_1\) & 0.7\% & 7.8\% \\
\(q_9\) & 6.7\% & 9.8\% \\
\(q_{10}\) & 3.2\% & 6.0\% \\
\bottomrule
\end{tabular}
\end{table}

We write \(q_N\) for the physical qubit addressed as \texttt{\$N} in
OpenQASM 3. The mean is about 6.5\% per qubit, roughly three times the
assumed value and asymmetric in the direction expected from \(T_1\) decay
during readout; feeding it back moves the predicted purity from 0.89 to
about 0.75, against the measured 0.7184. Readout is also correlated:
\(q_0\)'s \(P(1|0)\) rises from 1.6\% with idle neighbors to 16.9\% with
excited ones, and including that crosstalk predicts 0.7163 against the
measured 0.7184. The correlation saturates with excitation weight rather
than accumulating, so it must be measured at each weight; the
weight-resolved data is in the supplementary material. The input that was
wrong is precisely the one the specification does not publish.

\hypertarget{cross-session-drift-accounts-for-the-gap}{%
\subsection{Cross-session variation is the leading explanation for the
gap}\label{cross-session-drift-accounts-for-the-gap}}

With readout on the GHZ ladder measured in the same session rather than
assumed, we ran a bracketed experiment (Fig.~\ref{fig:hardware}): opening
calibration, predictions locked and committed from it, SWAP measurements,
closing calibration. Within-session drift measured about 1\% per qubit, so
the device was stable while the experiment ran. All three cells
nonetheless fell below both prediction bands, and the degradation was
non-monotonic: \(n = 3\) came out worse than \(n = 4\), despite the
\(n = 4\) register containing the \(n = 3\) one, an ordering neither
readout (which scales with the \(2n\) measured qubits) nor CZ count (which
grows with \(n\)) can produce. The band-by-band comparison is in the
supplementary material.

The cross-session picture explains it. Running byte-identical circuits on
the same physical qubits across three sessions, the \(n = 3\) register
healed (\(0.317 \to 0.378 \to 0.587\)) while \(n = 4\) drifted down
(\(0.431 \to 0.420 \to 0.382\)), and their ordering flipped. Within-session
drift propagates through the device model to about 0.01 in purity;
cross-session drift is approximately 0.2, roughly twenty times larger on
that like-for-like comparison. Identifying the non-monotonicity as a drift
artifact required running the same circuit three times over several days,
which is not standard practice.

\begin{figure}[tb]
\centering
\includegraphics[width=\linewidth]{fig6_hardware}
\caption{Hardware: cross-session variation is the leading explanation for the gap. (a)~Measured collective purity (vermillion, 95\% CI) at $n=2,3,4$ against the same-session prediction bands committed in advance from the opening (blue) and closing (orange) readout calibrations; every measurement falls below both bands. (b)~The same byte-identical circuits across three sessions: the $n=3$ register healed ($0.32\to0.59$) while $n=4$ drifted down ($0.43\to0.38$), and the ordering flipped. Cross-session drift ($\sim\!0.2$) is about $20\times$ the within-session drift (scale bar).}
\label{fig:hardware}
\end{figure}

The single-copy route failed the same way. A fifteen-basis \(n = 2\)
anchor, whose U-statistic our shadow pipeline predicted at
\(1.500 \pm 0.028\) from the same device parameters, returned \(1.344\)
with a 95\% bootstrap interval \([1.266, 1.433]\) over 9000 snapshots,
5.6 times the predicted spread, with disjoint intervals, and no simulated
run of the locked model reaches the measured value. That statistic is not
a purity, and the anchor tests this pipeline rather than the standard
i.i.d.\ estimator (supplementary material). We do not adjust the model to
fit. Two protocols with different circuits and estimators both missing
their locked predictions in one session is what makes a session-dependent
cause, rather than any protocol-specific defect, the parsimonious
explanation.

\hypertarget{the-elimination-ledger}{%
\subsection{The elimination
ledger}\label{the-elimination-ledger}}

Each candidate mechanism was addressed with a quantitative argument
rather than by absence of evidence: four are excluded outright, and three
are bounded well below the observed deficit.

\begin{longtable}[c]{@{}
  >{\raggedright\arraybackslash}p{0.5000\columnwidth}
  >{\raggedright\arraybackslash}p{0.5000\columnwidth}@{}}
\toprule
mechanism & how it was bounded or excluded \\
\midrule
\endfirsthead
\toprule
mechanism & how it was bounded or excluded \\
\midrule
\endhead
Within-session drift & Bracketed calibration measured approximately 1\%
per qubit, against a 0.2 discrepancy. \\
Single-copy readout error & Directly characterized. Feeding measured
values into the model does not close the gap at \(n \geq 3\). \\
Gate count & Cannot produce \(n=3\) worse than \(n=4\), since CZ count
is monotone in \(n\). \\
Cross-copy readout crosstalk & Ruled out by a physical bound. Joint
readout on the full \(2n\) register is factorizable (TVD 0.009 at
\(n=3\), 0.011 at \(n=4\)), and parity-pair correlations are negligible
and no stronger than non-pair correlations. Holding the measured
single-qubit flip rates fixed and pushing every copy-pair to its maximum
physically realizable correlation, which is bounded because
\(P(\text{both flip}) \leq \min\) of the two marginals and the measured
marginals are small, the entire achievable range of predicted purity is
\([0.611, 0.669]\) at \(n = 3\). The measurement is 0.378. Even maximal correlation leaves 0.23 of the deficit unexplained by this mechanism. \\
Coherent gate error & Randomized compiling. Pauli-twirled SWAP circuits,
with a positive control confirming that the inserted gates physically
execute (5.4:1 mass on the disjoint support), moved the purity by
\(-0.013\) against the 0.197 same-session \(n = 4\) gap, and twirl-to-twirl scatter was
shot-noise limited with zero physical scatter. A masking bound places
any hidden coherent contribution at no more than 14\% of the deficit.
The residual is incoherent. \\
A specific bad qubit pair (\(\{3, 12\}\)) & Localization test. An
\(n=3\) ladder that includes the suspect pair measured 0.581, matching
the standard \(n=3\) ladder's 0.581 to well within the \(\approx 0.012\)
shot-noise standard deviation at 5000 shots, from genuinely distinct raw
data. Adding the suspect qubits changes nothing. \\
A fixed register-geometry fault & The \(n=4\) versus \(n=3\) ordering
flips sign across sessions on identical registers. The deficit is
session-dependent, not geometric. \\
\bottomrule
\end{longtable}

\hypertarget{what-remains-and-what-we-conclude}{%
\subsection{What remains, and what we
conclude}\label{what-remains-and-what-we-conclude}}

What is left is a large, incoherent, session-dependent fault that appears
and disappears on the same physical qubits running byte-identical
circuits. We do not identify it; localizing it would require per-edge
interleaved randomized benchmarking across many sessions, a research
program rather than a section.

The conclusion the data support is narrower than ``we found a new error
mechanism'' and more useful than ``the device is noisy.'' Static
characterization did not predict collective-measurement performance over
the execution path available to us and across the sessions we ran. We
cannot isolate the mechanism: byte-identical circuits routed through a
decentralized network may have run on different operators' hardware, on
the same hardware in different calibration states, or both. The operational
consequence is the same either way: a protocol whose advantage rests on
a bounded, characterizable bias floor cannot be deployed against a floor
that moves by 0.2 between sessions, whatever moves it, and per-session
verification rather than published median parameters is what it requires.

Temporal and spatial instability of superconducting processors is
documented over 22 months \cite{dasgupta2021stability,dasgupta2022assessing}
and in calibration-drift analyses \cite{fewshot2026}, and we claim no
priority over it. What we add is a demonstration, with predictions
committed in advance and a seven-mechanism ledger, that a session-dependent
in this case study session-dependent variation dominated the modeled
collective-measurement error budget, at magnitudes that dwarf every
error source such protocols are conventionally modeled with, and
prevented a stable hardware test of the predicted advantage. The
single-copy anchor failed in the same session, so the variation is not
specific to the collective route.

One further constraint affects reproducibility rather than physics: the
single-copy route needs many independent measurement bases, and on a
per-circuit-priced platform without batching each is a separately billable
job, one to two orders of magnitude more than the collective route at
the same shot count. That is why our hardware single-copy baseline at
\(n \geq 3\) is computed from independently measured device parameters
rather than measured directly, with the fifteen-basis \(n = 2\) anchor the
one point we could afford. That anchor did not reproduce its prediction,
so the computed points are model output the hardware did not corroborate,
and we present them as such. The credit accounting is in the supplementary
material.

